% Generated draft: requires mathematical and source review.
\documentclass[11pt]{article}
\usepackage[margin=25mm]{geometry}
\usepackage{fontspec}
\setmainfont{Latin Modern Roman}
\usepackage{amsmath,amssymb,mathtools,unicode-math}
\setmathfont{Latin Modern Math}
\usepackage{enumitem,needspace,xurl,hyperref}
\setlength{\parindent}{0pt}
\setlength{\parskip}{5pt}
\setlength{\emergencystretch}{3em}
\newcommand{\N}{\mathbb{N}}
\newcommand{\Z}{\mathbb{Z}}
\newcommand{\Q}{\mathbb{Q}}
\newcommand{\R}{\mathbb{R}}
\newcommand{\C}{\mathbb{C}}
\newcommand{\F}{\mathbb{F}}
\newcommand{\A}{\mathbb{A}}
\newcommand{\PP}{\mathbb{P}}
\newcommand{\E}{\mathbb{E}}
\newcommand{\eps}{\varepsilon}
\newcommand{\dd}{\,\mathrm{d}}
\begin{document}
\section*{UnsolvedMath ID: 30004856}
\textbf{Source problem code:} catalogue entry
\subsection{Definitions and mathematical statement}
In the algebraic-geometric setting of the supplied record, a complex surface is a complex algebraic variety of dimension two. A nonempty surface \(S\) is \(\Q\)-acyclic if its complex points, with their analytic topology, satisfy
\[
H_0(S(\C);\Q)\cong\Q,\qquad H_i(S(\C);\Q)=0\quad\text{for every integer }i>0,
\]
where \(H_i\) denotes singular homology. Use the standard convention that a log-smooth completion \((X,D)\) consists of a smooth projective complex surface \(X\) and a reduced simple normal crossings divisor \(D\) with \(S\cong X\setminus\operatorname{Supp}(D)\). Thus the components of \(D\) are smooth and meet transversely, with at most two components through any point. Let \(K_X\) denote a canonical divisor of \(X\), and let \(\kappa\) denote the Iitaka dimension of a rational divisor. In particular,
\[
\kappa\!\left(K_X+\tfrac12D\right)=-\infty
\quad\Longleftrightarrow\quad
H^0\!\left(X,\mathcal O_X\bigl(m(2K_X+D)\bigr)\right)=0
\quad\text{for every integer }m\geq1.
\]
\textbf{Negativity Conjecture, as supplied.} For every complex \(\Q\)-acyclic surface \(S\) and every log-smooth completion \((X,D)\) of \(S\), is it true that
\[
\kappa\!\left(K_X+\tfrac12D\right)=-\infty?
\]
\textbf{Qualification concerning the domain.} Under the completion convention above, \(S\) is necessarily smooth because it is an open subset of the smooth surface \(X\). The supplied statement does not separately specify smoothness or define its completion convention. If singular \(\Q\)-acyclic surfaces are intended through a resolution-based meaning of completion, the required resolution and boundary conventions are missing; that extension is not formulated here.
\subsection{Short English statement}
The supplied Negativity Conjecture asks whether every complex \(\Q\)-acyclic surface with a log-smooth completion \((X,D)\) has \(\kappa(K_X+\tfrac12D)=-\infty\). Equivalently, every positive integral multiple of \(2K_X+D\) has no nonzero global sections. The standard completion convention forces the surface to be smooth; any intended singular-surface formulation requires clarification.
\subsection{Sources}
\begin{itemize}
\item Local UnsolvedMath registry, supplied source record for permanent ID \(30004856\), source problem code \texttt{catalogue entry}, titled ``Negativity of Log Divisors on Q-Acyclic Surfaces.'' The record supplies status ``open'' and statement status \texttt{corrected\_verified}, together with a claim of manual semantic, typographic, and LaTeX review. These are supplied registry assessments, not independent verification here.
\item Supplied report citation: \emph{Singularities (2022)}, Oberwolfach Reports, DOI 10.4171/owr/2021/47, pp. 2573--2639, \url{https://doi.org/10.4171/owr/2021/47}. The record describes Oberwolfach Reports as documenting workshops and their open-problem discussions. It supplies proposed year 2021 and no named proposer. Both the cited 2022 title and the supplied 2021 year are retained without reconciliation.
\item The report excerpt labels the conjecture ``Negativity Conjecture, [Pal19, 4.7].'' No full bibliographic entry for [Pal19] is supplied, so its author and publication details cannot be identified from this record alone.
\item Supplied literature source: ``Pełka, smooth Q-homology planes satisfying negativity,'' \url{https://doi.org/10.1112/jlms.12779}. The supplied assessment says that Pełka classifies the smooth \(\Q\)-homology planes satisfying the conjecture; this is a conditional classification, not a proof of the universal assertion. The excerpt attributes the corresponding thesis work to T. Pełka and mentions almost log exceptional curves as a tool for grouping such surfaces and rational cuspidal curves into series.
\item Supplied literature source: ``OWR statement that the conjecture remains open,'' \url{https://doi.org/10.4171/owr/2021/47}.
\item Literature assessment date: \texttt{2026-08-22}. The supplied dated curation describes the conjecture as still open and distinguishes Pełka's conditional classification from the universal conjecture. This assessment is supplied dated curation, not independently verified; the record explicitly requests independent verification of time-sensitive status and cited claims.
\end{itemize}
\end{document}
